% Openers (design 5e). Part pages and chapter openers share this file.
%
% Part page: a full charcoal recto. Masthead, volume line, the coral part numeral,
% PART n, the ivory title and the part rule. No epigraph, and no illustration: the
% cover motif and \chaptermotif are not drawn here.
%
% Chapter opener: the same lockup, compact, at the top of the chapter's first page.
% A charcoal band (masthead, volume line, coral numeral) sits 0.30 in inside the trim,
% then CHAPTER n, the Cormorant title, a coral rule and the epigraph. The chapter
% prose starts immediately under that block, on the same page. There is no second
% heading. The first page uses the plain style (folio at the foot, no running head).
% Numbered chapters start on the next page, recto or verso: the class option is
% openany, and a blank is not inserted to force a side. A part's first chapter
% therefore still sits on the back of the part page.
%
% Epigraph: the chapter's opening \begin{center}\itshape ... \end{center} (directly after
% \chapter{...} and its optional \label) is set under the rule; the source is unchanged.
%
% \chaptermotif is accepted and ignored. The print openers do not draw it, and they
% do not input cover-motif.tex.
\makeatletter
\newlength{\oc@opL}
\newsavebox{\oc@pt@text}
% Chapter band. The ink stays 0.30 in inside the trim on the top and the sides, so a
% no-bleed KDP page does not carry the band to the cut. \oc@bandH is the band's height;
% \oc@bandDrop is how far its top sits below the trim (the text block begins at 1 in).
\newlength{\oc@bandInset}\setlength{\oc@bandInset}{0.30in}
\newlength{\oc@bandH}\setlength{\oc@bandH}{1.48in}
\newlength{\oc@bandDrop}\setlength{\oc@bandDrop}{0.30in}
% How far titlesec's strut line and the 0 pt before-skip push the band down
% past \oc@bandDrop. Measured on a 7 x 10 page (see the comment on \oc@flow@band).
\newlength{\oc@bandRaise}\setlength{\oc@bandRaise}{2.05in}
\newlength{\oc@bandW}
\newlength{\oc@hang}
% Long titles: allow a line break after an en or em dash ("Bombieri--Vinogradov").
\ExplSyntaxOn
\tl_new:N \l__oc_op_title_tl
\cs_new_protected:Npn \ocBreakDashes #1
  {
    \tl_set:Nn \l__oc_op_title_tl {#1}
    \tl_replace_all:Nnn \l__oc_op_title_tl { --- } { \textemdash\allowbreak }
    \tl_replace_all:Nnn \l__oc_op_title_tl { -- } { \textendash\allowbreak }
    \tl_use:N \l__oc_op_title_tl
  }
\ExplSyntaxOff
% The print opener carries no motif and no volume illustration.
\newcommand{\chaptermotif}[1]{}
\newif\ifoc@op@recto
\let\oc@ch@epi\@empty
% ---------------------------------------------------------------- the chapter band
% Drawn in the text flow and shifted into the top and side margins, stopping
% \oc@bandDrop / \oc@bandInset short of the trim. The following title block, and
% then the chapter prose, continue on this page.
\newcommand{\oc@flow@band}{%
  \setlength{\oc@bandW}{\dimexpr\PRINTWIDTH in-2\oc@bandInset\relax}%
  \ifodd\c@page
    \setlength{\oc@hang}{\dimexpr\ocInner-\oc@bandInset\relax}%
  \else
    \setlength{\oc@hang}{\dimexpr\ocOuter-\oc@bandInset\relax}%
  \fi
  % titlesec inserts \vspace*{0pt} and its block shape leaves a strut line above
  % this point, so the natural top of the band is far down the page. \oc@bandRaise
  % is the measured correction that puts the band's top at \oc@bandDrop.
  \vspace*{\dimexpr\oc@bandH-(1in-\oc@bandDrop)+\topskip-\oc@bandRaise\relax}%
  \nointerlineskip
  \moveleft\oc@hang\hbox{%
    \begin{tikzpicture}[baseline={(0,0)}]
      \useasboundingbox (0,0) rectangle (\oc@bandW,\oc@bandH);
      \fill[ocBlack] (0,0) rectangle (\oc@bandW,\oc@bandH);
      \node[anchor=north west,inner sep=0pt] at (\oc@hang,\oc@bandH-0.13in)
        {\ocDisplay\fontsize{10.5}{12}\selectfont\color{ocCream}\ocCaps[26]{Order and Chance}};
      \node[anchor=north west,inner sep=0pt] at (\oc@hang,\oc@bandH-0.34in)
        {\ocBook\fontsize{7}{9}\selectfont\color{ocCream!80}\ocCaps[20]{\ocVolume}};
      \node[anchor=base west,inner sep=0pt,text=ocCoral,font=\ocNumeral\fontsize{78}{78}\selectfont]
        at (\oc@hang-0.015in,0.09in) {\thechapter};
    \end{tikzpicture}}}
% Label, title, coral rule and the epigraph, on the paper directly under the band.
\newcommand{\oc@flow@title}[1]{%
  \nointerlineskip
  \vskip 11pt
  \vbox{%
    \hsize=\linewidth
    \everypar{}%
    \parindent=0pt
    \noindent{\ocLabelFont\ocLab{\chaptertitlename}\enspace\ocNum{\thechapter}}\par
    \vspace{6pt}%
    {\ocDisplay\linespread{1}\fontsize{30}{34}\selectfont\color{ocBlack}%
      \raggedright\hyphenpenalty10000\exhyphenpenalty10000 \ocBreakDashes{#1}\par}%
    \vspace{8pt}%
    \ocChapterRule\par
    \ifx\oc@ch@epi\@empty\else
      \vspace{8pt}%
      {\small\itshape\color{ocBlack}\raggedright\hyphenpenalty10000\exhyphenpenalty10000
        \let\\\newline \oc@ch@epi\par}%
    \fi}}
% Numbered chapter heading: plain page (no running head), then the band and title.
% \ocChapterHead is the titlesec before-code; #1 is the title that follows it.
\renewcommand{\ocChapterHead}[1]{%
  \thispagestyle{plain}%
  % titlesec's block shape has already \leavevmode'd, so the band has to return
  % to vertical mode. The paragraph it closes holds only the indent strut.
  \par
  \oc@flow@band
  \oc@flow@title{#1}}
% Part page: full charcoal, the part's own title lockup, nothing else.
\newcommand{\oc@part@page}[2]{% #1 numeral (empty for \part*), #2 title
  \global\oc@op@rectotrue
  \global\setbox\oc@pt@text\hbox{\parbox[t]{4.3in}{%
    \if\relax\detokenize{#1}\relax\else
    {\ocMed\fontsize{8.2}{10}\selectfont\color{ocCream}\ocLab{\partname}\enspace\textcolor{ocCoral}{#1}}\par
    \vspace{10pt}\fi
    {\ocDisplay\linespread{1}\fontsize{38}{42}\selectfont\color{ocCream}\raggedright\hyphenpenalty10000\exhyphenpenalty10000 \ocBreakDashes{#2}\par}%
    \vspace{16pt}\oc@op@diarule{ocCoral}\par}}%
  \xdef\oc@op@numeral{#1}%
  \AddToHookNext{shipout/background}{\put(0,0){\oc@part@art}}%
  \thispagestyle{empty}\null\clearpage}
\newcommand{\oc@part@art}{%
  \ifoc@op@recto\setlength{\oc@opL}{\ocInner}\else\setlength{\oc@opL}{\ocOuter}\fi
  \begin{tikzpicture}[baseline={(0,0)}]
    \useasboundingbox (0,0) rectangle (\paperwidth,-\paperheight);
    \fill[ocBlack] (0,0) rectangle (\paperwidth,-\paperheight);
    \begin{scope}[shift={(\ocBleed,-\ocBleed)}]
      \node[anchor=north west,inner sep=0pt] at (\oc@opL,-0.75in)
        {\ocDisplay\fontsize{10.5}{12}\selectfont\color{ocCream}\ocCaps[26]{Order and Chance}};
      \node[anchor=north west,inner sep=0pt] at (\oc@opL,-0.98in)
        {\ocBook\fontsize{7}{9}\selectfont\color{ocCream!80}\ocCaps[20]{\ocVolume}};
      \node[anchor=base west,inner sep=0pt,text=ocCoral,font=\ocNumeral\fontsize{190}{190}\selectfont]
        at (\oc@opL-0.08in,-3.35in) {\oc@op@numeral};
      \node[anchor=north west,inner sep=0pt] at (\oc@opL,-3.75in) {\usebox\oc@pt@text};
    \end{scope}
  \end{tikzpicture}}
% The part page's short rule: rule, diamond, rule.
\newcommand{\oc@op@diarule}[1]{\tikz[baseline=-0.6ex]{\draw[#1,line width=1.2pt](0,0)--(0.36in,0);
  \fill[#1](0.45in,0.045in)--(0.495in,0)--(0.45in,-0.045in)--(0.405in,0)--cycle;\draw[#1,line width=1.2pt](0.54in,0)--(0.9in,0);}}
% ---------------------------------------------------------------- \chapter wrapper
% Reads \chapter[short]{title}, an optional \label and an opening epigraph, typesets the
% opener (numbered main-matter chapters only), then calls the original \chapter.
\def\oc@centername{center}
\def\oc@skip#1{\def\oc@skip@k{#1}\oc@skip@loop}
\def\oc@skip@loop{\@ifnextchar\par{\oc@skip@par}{\oc@skip@k}}
\long\def\oc@skip@par#1{\oc@skip@loop}
\def\oc@ch@go{\@ifnextchar[{\oc@ch@opt}{\oc@ch@noopt}}
\long\def\oc@ch@noopt#1{\oc@ch@a{#1}{#1}}
\long\def\oc@ch@opt[#1]#2{\oc@ch@a{#1}{#2}}
\long\def\oc@ch@a#1#2{\def\oc@ch@short{#1}\def\oc@ch@title{#2}%
  \let\oc@ch@label\@empty\let\oc@ch@epi\@empty\oc@skip\oc@ch@lab}
\def\oc@ch@lab{\@ifnextchar\label{\oc@ch@getlab}{\oc@ch@afterlab}}
\def\oc@ch@getlab\label#1{\def\oc@ch@label{\label{#1}}\oc@skip\oc@ch@afterlab}
\def\oc@ch@afterlab{\@ifnextchar\begin{\oc@ch@begin}{\oc@ch@finish{}}}
\def\oc@ch@begin\begin#1{\def\oc@tmp{#1}%
  \ifx\oc@tmp\oc@centername\expandafter\@firstoftwo\else\expandafter\@secondoftwo\fi
  {\@ifnextchar\itshape{\oc@ch@epigraph}{\oc@ch@finish{\begin{center}}}}%
  {\oc@ch@finish{\begin{#1}}}}
\long\def\oc@ch@epigraph\itshape#1\end#2{\def\oc@tmp{#2}%
  \ifx\oc@tmp\oc@centername\expandafter\@firstoftwo\else\expandafter\@secondoftwo\fi
  {\def\oc@ch@epi{#1}\oc@ch@finish{}}%
  {\oc@ch@finish{\begin{center}\itshape#1\end{#2}}}}
\long\def\oc@ch@finish#1{%
  \if@mainmatter
    % Next page, either side. \chapter would otherwise honour \@openright and
    % insert a blank verso. Restore the flag before the body is read, so the
    % heading's \everypar (no indent on the first paragraph) survives.
    \@openrightfalse
    \expandafter\oc@chapter@orig\expandafter[\expandafter{\oc@ch@short}]{\oc@ch@title}%
    \@openrighttrue
    \oc@ch@label#1%
  \else
    \expandafter\oc@chapter@orig\expandafter[\expandafter{\oc@ch@short}]{\oc@ch@title}\oc@ch@label
    \ifx\oc@ch@epi\@empty\else\begin{center}\itshape\oc@ch@epi\end{center}\fi#1%
  \fi}
\AtBeginDocument{\let\oc@chapter@orig\chapter
  \renewcommand{\chapter}{\@ifstar{\oc@chapter@orig*}{\oc@ch@go}}}
\makeatother
